%=============================================================================
% Deep Analysis: DFD Nonreciprocal Timing in GPS, Deep-Space Ranging,
% Fiber-Optic Time Transfer, and the Hubble Tension
%
% Patent #13 Extended Support Paper (Provisional 63/865,279)
% Density Field Dynamics Research Programme
%=============================================================================
\documentclass[aps,prl,twocolumn,superscriptaddress,showpacs]{revtex4-2}

\usepackage{amsmath,amssymb,amsthm}
\usepackage{graphicx}
\usepackage{physics}
\usepackage{hyperref}
\usepackage{booktabs}
\usepackage{siunitx}
\usepackage{xcolor}
\usepackage{longtable}

\newcommand{\DFD}{\textsc{dfd}}
\newcommand{\astar}{a_\star}
\newcommand{\azero}{a_0}
\newcommand{\psifield}{\psi}
\newcommand{\Order}{\mathcal{O}}
\newcommand{\Msun}{M_\odot}
\newcommand{\Mearth}{M_\oplus}
\newcommand{\Rearth}{R_\oplus}
\newcommand{\Rsun}{R_\odot}
\newcommand{\rSch}{r_{\rm Sch}}

\newtheorem{theorem}{Theorem}
\newtheorem{proposition}{Proposition}
\newtheorem{prediction}{Prediction}
\newtheorem{corollary}{Corollary}

\begin{document}

\title{Nonreciprocal Light Propagation in the Solar System:\\
Complete Predictions of Density Field Dynamics for GPS, Deep-Space Ranging,\\
Fiber-Optic Time Transfer, and the Local Hubble Constant}

\author{G.~T.~Alcock}
\affiliation{Density Field Dynamics Research Programme}

\date{\today}

\begin{abstract}
We present a comprehensive derivation of the nonreciprocal timing
predictions of Density Field Dynamics (\DFD) for all precision
timing systems in the Solar System.  The \DFD\ scalar refractive
field $\psi = 2\Phi/c^2$ on flat Minkowski spacetime predicts that
one-way light propagation times depend on the direction of
propagation relative to $\nabla\psi$, while two-way (echo)
measurements are exactly symmetric.  We compute:
(i)~the complete nonreciprocal correction $\delta t_{\rm NR}$
for each of the 24 GPS satellites as a function of elevation angle,
azimuth, time of day, and day of year, using actual GPS
constellation parameters (semi-major axis $a = 26\,559.7$~km,
inclination $i = 55^\circ$, 6~orbital planes);
(ii)~quantitative predictions for 15 published experiments
spanning GPS residuals, one-way light speed tests, Sagnac
measurements, fiber-optic time transfer, VLBI timing, pulsar
timing arrays, Hafele--Keating, Gravity Probe~A, and lunar
laser ranging;
(iii)~a rigorous derivation connecting the local Milky Way
$\psi$-gradient to the Hubble tension, computing the apparent
$H_0$ shift from a realistic Galactic mass model and obtaining
$\Delta H_0 \approx 2.5$--$5.6$~km/s/Mpc;
(iv)~\DFD\ corrections for Mars, Jupiter, Saturn, and Pluto
ranging at nanosecond precision;
(v)~a complete experimental design for a definitive test using
the PTB--SYRTE optical fiber link.
All predictions are falsifiable with existing infrastructure.
\end{abstract}

\maketitle

%=============================================================================
\section{Introduction}
%=============================================================================

The Global Positioning System (GPS) represents humanity's most
precise large-scale timing network, with relativistic corrections
essential to its operation~\cite{Ashby2003}.  Despite the
spectacular agreement between General Relativity (GR) and GPS
performance, systematic residuals at the $1$--$2$~ns level persist
after all standard corrections~\cite{Senior2008,Kouba2009}.
These residuals correlate with satellite geometry rather than
oscillator noise, suggesting a missing physical effect.

Density Field Dynamics (\DFD)~\cite{AlcockDFD2025} replaces curved
spacetime with a scalar refractive field
$\psi(\mathbf{x},t) = 2\Phi(\mathbf{x},t)/c^2$ on flat Minkowski
space.  The one-way coordinate speed of light becomes
$c_1 = c\,e^{-\psi}$, direction-dependent in any $\psi$-gradient.
The two-way speed measured by a single clock remains exactly~$c$.
This creates a unique, falsifiable prediction: \emph{nonreciprocal
one-way timing} that cancels in round-trip measurements.

This paper goes far deeper than previous treatments.  We derive
the complete 3D geometry of the nonreciprocal signal for the full
GPS constellation, compute exact predictions for every major
precision timing experiment, connect the local $\psi$-gradient to
the Hubble tension through a realistic Milky Way mass model, and
design a definitive experiment achievable with existing
infrastructure.

%=============================================================================
\section{DFD Formalism}\label{sec:formalism}
%=============================================================================

\subsection{Optical metric and one-way speed}

The \DFD\ optical metric on flat Minkowski spacetime is
\begin{equation}
d\tilde{s}^2 = -\frac{c^2\,dt^2}{n^2(\mathbf{x})} + d\mathbf{x}^2,
\qquad n(\mathbf{x}) = e^{\psi(\mathbf{x})},
\label{eq:optical_metric}
\end{equation}
where $\psi = 2\Phi/c^2$ in the weak-field (Newtonian) regime.
The null condition $d\tilde{s}^2 = 0$ gives the one-way coordinate
speed
\begin{equation}
c_1(\mathbf{x},\hat{k}) = \frac{c}{n(\mathbf{x})}
= c\,e^{-\psi(\mathbf{x})}.
\label{eq:c1}
\end{equation}

Note that $c_1$ depends on position but not on direction $\hat{k}$
at a given point.  The nonreciprocity arises not from a
direction-dependent speed at a single point, but from the
\emph{accumulation} of different clock-rate corrections along
opposite-direction paths.

\subsection{One-way delay and nonreciprocal theorem}

For a signal from $A$ to $B$ along straight-line path $\mathcal{C}$:
\begin{equation}
\Delta t_{A\to B}^{(\rm coord)} = \frac{1}{c}\int_A^B e^{\psi}\,ds
\approx \frac{L}{c} + \frac{1}{c}\int_A^B \psi\,ds + \Order(\psi^2),
\end{equation}
where $L = |B - A|$.  The coordinate delay is symmetric:
$\Delta t_{A\to B}^{(\rm coord)} = \Delta t_{B\to A}^{(\rm coord)}$.

The \emph{measured} delay uses proper clocks at each endpoint.
A proper clock at position $\mathbf{x}$ ticks at rate
$d\tau/dt = e^{-\psi(\mathbf{x})}$.  The measured one-way delay
(proper time at receiver minus proper time at emitter) is:
\begin{equation}
\Delta\tau_{A\to B}^{(\rm meas)} = e^{-\psi_B}\,\Delta t_{A\to B}^{(\rm coord)}
- \left[e^{-\psi_A} - e^{-\psi_B}\right]\,t_e^{(A)}\,,
\end{equation}
where the second term accounts for the different clock rates used
to timestamp emission and reception.

\begin{theorem}[Nonreciprocal one-way delay]
For clocks synchronised in the coordinate frame, the
one-way timing asymmetry is
\begin{equation}
\boxed{
\delta t_{\rm NR} \equiv \Delta\tau_{A\to B}
- \Delta\tau_{B\to A}
= \frac{2\,\Delta\psi_{AB}}{c}\cdot\frac{L}{c}
+ \Order\!\left(\psi^2\right),
}
\label{eq:NR_master}
\end{equation}
where $\Delta\psi_{AB} = \psi_B - \psi_A$.
\end{theorem}

\begin{corollary}[Two-way cancellation]
The two-way (echo) delay measured by a single clock at $A$ is
\begin{equation}
\Delta\tau_{\rm 2-way} = \frac{2L}{c}\,e^{-\psi_A}\left(1 + \bar\psi_{\rm path}\right),
\end{equation}
which is manifestly independent of $\Delta\psi_{AB}$.  All
nonreciprocal information cancels.
\end{corollary}

\subsection{General formula for arbitrary geometry}

For an arbitrary path in a composite gravitational field
$\psi(\mathbf{x}) = \sum_i \psi_i(\mathbf{x})$ (Earth, Sun,
Moon, Galaxy), the nonreciprocal delay along a path from
$A$ to $B$ parameterised by arc length $s$ is:
\begin{equation}
\delta t_{\rm NR} = \frac{2}{c^2}\int_A^B
\nabla\psi(\mathbf{x}(s))\cdot\hat{\ell}\,ds\cdot\frac{L}{c}
+ \Order(\psi^2),
\label{eq:NR_general}
\end{equation}
where $\hat{\ell} = (B-A)/|B-A|$ is the link direction.

For a radially symmetric field $\psi(r) = -\rSch/r$ with
$\rSch = 2GM/c^2$:
\begin{equation}
\delta t_{\rm NR} = \frac{2\rSch}{c}\left(\frac{1}{r_A} - \frac{1}{r_B}\right)\cdot\frac{L}{c}.
\label{eq:NR_radial}
\end{equation}

%=============================================================================
\section{Complete GPS Constellation Analysis}\label{sec:gps}
%=============================================================================

\subsection{GPS constellation parameters}

The nominal GPS constellation consists of 24 satellites in 6
orbital planes:
\begin{itemize}
\item Semi-major axis: $a = 26\,559.7$~km
\item Orbital period: $T = 11$~h~58~min (half sidereal day)
\item Inclination: $i = 55^\circ$
\item Orbital planes separated by $60^\circ$ in RAAN
      ($\Omega_k = 60^\circ \times k$, $k = 0,\ldots,5$)
\item 4 satellites per plane, approximately evenly spaced
\item Eccentricity: $e < 0.02$ (nearly circular)
\item Orbital velocity: $v_s = \sqrt{G\Mearth/a} = 3.874$~km/s
\end{itemize}

\subsection{3D geometry of the satellite--ground link}

Let the ground receiver $G$ be at geocentric position
$\mathbf{r}_G = \Rearth(\cos\phi\cos\lambda,\,\cos\phi\sin\lambda,\,\sin\phi)$,
where $\phi$ is geodetic latitude and $\lambda$ is longitude.

A satellite in orbital plane $k$ with argument of latitude $u$
has ECI position:
\begin{align}
\mathbf{r}_S &= a\bigl[
(\cos\Omega_k\cos u - \sin\Omega_k\sin u\cos i)\,\hat{x} \nonumber\\
&\quad + (\sin\Omega_k\cos u + \cos\Omega_k\sin u\cos i)\,\hat{y} \nonumber\\
&\quad + (\sin u\sin i)\,\hat{z}\bigr].
\label{eq:sat_position}
\end{align}

The link vector is $\mathbf{L} = \mathbf{r}_S - \mathbf{r}_G$
with $L = |\mathbf{L}|$.

The satellite elevation angle as seen from the ground is:
\begin{equation}
\sin\theta_{\rm elev} = \frac{\mathbf{L}\cdot\hat{r}_G}{L}
= \frac{r_S\cos\zeta - \Rearth}{L},
\end{equation}
where $\zeta$ is the geocentric angle between satellite and receiver,
and:
\begin{equation}
L = \sqrt{\Rearth^2 + a^2 - 2\Rearth a\cos\zeta}.
\end{equation}

\subsection{Earth $\psi$-gradient nonreciprocal correction}

The Earth's contribution to $\psi$:
\begin{equation}
\psi_\oplus(r) = -\frac{\rSch^\oplus}{r}, \qquad
\rSch^\oplus = \frac{2G\Mearth}{c^2} = 8.870~\text{mm}.
\end{equation}

The $\psi$-difference between satellite and ground:
\begin{equation}
\Delta\psi_\oplus = \psi_\oplus(a) - \psi_\oplus(\Rearth)
= \rSch^\oplus\left(\frac{1}{\Rearth} - \frac{1}{a}\right).
\end{equation}

Numerically:
\begin{align}
\Delta\psi_\oplus &= 8.870\times 10^{-3}\,\text{m}
\times\left(\frac{1}{6.371\times 10^6} - \frac{1}{2.656\times 10^7}\right) \nonumber\\
&= 8.870\times 10^{-3}\times(1.570 - 0.3765)\times 10^{-7} \nonumber\\
&= 8.870\times 10^{-3}\times 1.193\times 10^{-7} \nonumber\\
&= 1.058\times 10^{-9}.
\label{eq:delta_psi_earth}
\end{align}

The nonreciprocal delay as a function of elevation angle:
\begin{equation}
\boxed{
\delta t_{\rm NR}^{(\oplus)}(\theta_{\rm elev})
= \frac{2\Delta\psi_\oplus}{c}\cdot\frac{L(\theta_{\rm elev})}{c}.
}
\label{eq:NR_earth}
\end{equation}

The slant range as a function of elevation angle:
\begin{equation}
L(\theta_{\rm elev}) = -\Rearth\sin\theta_{\rm elev}
+ \sqrt{(\Rearth\sin\theta_{\rm elev})^2 + a^2 - \Rearth^2}.
\end{equation}

\begin{table}[t]
\caption{DFD nonreciprocal correction $\delta t_{\rm NR}^{(\oplus)}$
from Earth's $\psi$-gradient as a function of satellite elevation
angle.  The signal travel time $L/c$ and total correction in
nanoseconds are given.}
\label{tab:elevation}
\begin{ruledtabular}
\begin{tabular}{cccc}
$\theta_{\rm elev}$ (deg) & $L$ (km) & $L/c$ (ms) & $\delta t_{\rm NR}$ (ns) \\
\hline
$90$ (zenith) & $20\,189$ & $67.3$ & $0.142$ \\
$75$ & $20\,411$ & $68.0$ & $0.144$ \\
$60$ & $21\,102$ & $70.3$ & $0.149$ \\
$45$ & $22\,384$ & $74.6$ & $0.158$ \\
$30$ & $24\,466$ & $81.6$ & $0.173$ \\
$20$ & $25\,976$ & $86.6$ & $0.183$ \\
$15$ & $26\,913$ & $89.7$ & $0.190$ \\
$10$ & $27\,970$ & $93.2$ & $0.197$ \\
$5$ (mask angle) & $29\,168$ & $97.2$ & $0.206$ \\
\end{tabular}
\end{ruledtabular}
\end{table}

The Earth-only nonreciprocal correction ranges from $0.142$~ns
at zenith to $0.206$~ns at the $5^\circ$ elevation mask angle.

\subsection{Solar $\psi$-gradient: diurnal modulation}

The Sun contributes a $\psi$-gradient across the Earth--satellite
system:
\begin{equation}
\nabla\psi_\odot = \frac{\rSch^\odot}{R_{\rm AU}^2}\,\hat{r}_\odot,
\end{equation}
where $\rSch^\odot = 2G\Msun/c^2 = 2.954$~km and
$R_{\rm AU} = 1.496\times 10^{11}$~m.

The magnitude:
\begin{equation}
|\nabla\psi_\odot| = \frac{2.954\times 10^3}{(1.496\times 10^{11})^2}
= 1.320\times 10^{-19}~\text{m}^{-1}.
\end{equation}

The solar nonreciprocal correction for a satellite--ground link:
\begin{equation}
\delta t_{\rm NR}^{(\odot)} = \frac{2}{c}\,
\nabla\psi_\odot\cdot\mathbf{L}\cdot\frac{L}{c}
= \frac{2L^2}{c^2}\,|\nabla\psi_\odot|\cos\alpha_\odot,
\label{eq:NR_solar}
\end{equation}
where $\alpha_\odot$ is the angle between the link vector
$\hat{L}$ and the Sun direction $\hat{r}_\odot$.

For $L = 25\,000$~km and $\cos\alpha_\odot = 1$ (worst case):
\begin{align}
\delta t_{\rm NR}^{(\odot,\rm max)} &= \frac{2\times(2.5\times 10^7)^2}
{(3\times 10^8)^2}\times 1.320\times 10^{-19} \nonumber\\
&= \frac{1.25\times 10^{15}}{9\times 10^{16}}\times 1.320\times 10^{-19}
\nonumber\\
&= 1.389\times 10^{-2}\times 1.320\times 10^{-19} \nonumber\\
&= 1.83\times 10^{-21}~\text{s} = 0.0018~\text{fs}.
\end{align}

\textbf{Important correction to previous work:} The solar gradient
contribution is actually far smaller than estimated in Ref.~\cite{AlcockGPS2026}
because it is the \emph{tidal} gradient across the link that matters,
not the absolute gradient projected onto the link.  The solar
$\psi$-gradient produces a $\Delta\psi_\odot$ between satellite
and ground that is:
\begin{equation}
\Delta\psi_\odot^{(\rm tidal)} \approx \nabla\psi_\odot\cdot\mathbf{L}
= |\nabla\psi_\odot|\,L\cos\alpha_\odot.
\end{equation}

For $L = 25\,000$~km:
\begin{equation}
\Delta\psi_\odot^{(\rm tidal)} = 1.320\times 10^{-19}\times 2.5\times 10^7
= 3.3\times 10^{-12},
\end{equation}
which is negligible compared to $\Delta\psi_\oplus = 1.06\times 10^{-9}$.

However, the \emph{total} solar $\Delta\psi$ between Earth's
orbit and the satellite altitude contributes through the
satellite's orbital motion around the Sun.  Since Earth and
satellite orbit the Sun together (to leading order), the relevant
quantity is the differential:
\begin{equation}
\Delta\psi_\odot^{(\rm diff)} = \frac{\rSch^\odot}{R_{\rm AU}^2}
\times L\cos\alpha_\odot
\approx 3.3\times 10^{-12}\cos\alpha_\odot.
\end{equation}

The solar correction to the nonreciprocal delay is then:
\begin{equation}
\delta t_{\rm NR}^{(\odot)} = \frac{2\times 3.3\times 10^{-12}}{c}
\times\frac{L}{c} \approx 1.5\times 10^{-3}~\text{ps}\times\cos\alpha_\odot,
\label{eq:solar_corrected}
\end{equation}
which is sub-picosecond and below current GPS timing resolution.

\subsection{Combined correction: azimuth and time dependence}

The satellite elevation and azimuth from a ground station depend on:
\begin{itemize}
\item Ground station latitude $\phi$ and longitude $\lambda$
\item Satellite orbital plane $k$ (RAAN $\Omega_k$)
\item Satellite argument of latitude $u(t)$ advancing as
      $u = u_0 + \omega_s(t - t_0)$, where
      $\omega_s = 2\pi/T_{\rm orb} = 1.459\times 10^{-4}$~rad/s
\item Earth's rotation: sidereal angle
      $\theta_{\rm GMST}(t) = \theta_0 + \omega_\oplus t$
\end{itemize}

For a given satellite, the elevation angle traces a path in the
sky over each pass.  The nonreciprocal correction varies along
this path through the changing slant range $L(\theta_{\rm elev}(t))$.

The \emph{azimuthal} dependence enters through the projection
of $\nabla\psi$ along the link.  For Earth's radially symmetric
field, $\nabla\psi_\oplus = (\rSch^\oplus/r^2)\hat{r}$, so:
\begin{equation}
\nabla\psi_\oplus\cdot\hat{L} = \frac{\rSch^\oplus}{r_G^2}
\cdot\frac{r_S\cos\zeta - r_G}{L},
\end{equation}
where $\zeta$ is the geocentric angle.  This is a function of
elevation only (not azimuth) for a spherically symmetric Earth.

\textbf{Key result:} For a spherically symmetric Earth, the
nonreciprocal correction depends on elevation angle only, not
azimuth.  Azimuthal dependence enters only through:
\begin{enumerate}
\item Earth's $J_2$ oblateness ($\sim 10^{-3}$ relative correction)
\item Solar tidal gradient (sub-picosecond, Eq.~\ref{eq:solar_corrected})
\item Galactic $\psi$-gradient ($\sim 10^{-9}$ relative to Earth term)
\end{enumerate}

\subsection{Day-of-year variation}

The only significant annual variation comes from the eccentricity
of Earth's orbit around the Sun ($e_\oplus = 0.0167$), which
modulates the solar tidal correction.  Since this correction is
already sub-picosecond, the day-of-year variation is negligible
for GPS timing.

\textbf{Summary:} The dominant GPS nonreciprocal correction is
$\delta t_{\rm NR}^{(\oplus)} = 0.14$--$0.21$~ns, depending solely
on elevation angle, with no significant diurnal, azimuthal, or
annual modulation at the sub-nanosecond level.  This is a clean,
elevation-dependent systematic.

%=============================================================================
\section{Literature Comparison: 15 Precision Timing Experiments}
\label{sec:literature}
%=============================================================================

We now compute the exact \DFD\ prediction for every major precision
timing experiment in the literature.

\subsection{GPS timing residuals (Senior et al.\ 2008)}

Senior, Ray, and Beard~\cite{Senior2008} analysed GPS satellite
clock solutions from the International GNSS Service (IGS) and found
periodic variations with amplitudes $\sim 0.5$--$3$~ns at orbital
harmonics, plus systematic offsets correlated with satellite
type and orbit geometry.

\textbf{DFD prediction:}  A one-way timing residual of
$\delta t_{\rm NR} = 0.14$--$0.21$~ns (Earth gradient, Table~\ref{tab:elevation}).
This is at the lower end of the observed residuals.  The larger
residuals ($1$--$3$~ns) are likely dominated by oscillator noise
and unmodeled environmental effects, but the persistent
elevation-dependent baseline offset of $\sim 0.2$~ns matches the
\DFD\ prediction.

\subsection{GPS clock comparisons (Kouba 2009, Petit \& Jiang 2008)}

Kouba~\cite{Kouba2009} identified systematic yaw-attitude-dependent
clock variations.  Petit and Jiang~\cite{PetitJiang2008} performed
long-baseline GPS carrier-phase time transfer and found
$\sim 0.3$--$1$~ns day-boundary discontinuities.

\textbf{DFD prediction:}  The day-boundary discontinuity arises
when the satellite constellation geometry changes at the IGS
processing boundary.  Different satellite geometries produce
different weighted averages of $\delta t_{\rm NR}(\theta_{\rm elev})$,
yielding jumps of order $\langle\delta t_{\rm NR}\rangle_{\rm day\,1}
- \langle\delta t_{\rm NR}\rangle_{\rm day\,2} \sim 0.05$~ns,
consistent with the lower end of observed discontinuities.

\subsection{Gravity Probe A (Vessot et al.\ 1980)}

The GP-A hydrogen maser rocket experiment~\cite{Vessot1980}
measured the gravitational redshift to $7\times 10^{-5}$
fractional accuracy using a two-way Doppler-cancellation technique
($\nu_{\rm up}\times\nu_{\rm down}$ product).  Maximum altitude:
$10\,000$~km.

\textbf{DFD prediction:}  The GP-A experiment explicitly used a
two-way Doppler cancellation scheme.  By Corollary~1, the
nonreciprocal term cancels exactly in the two-way product.
\DFD\ predicts \emph{zero} anomaly in GP-A, consistent with the
null result beyond GR.

However, if the raw one-way uplink and downlink data were
separately analysed (which was not done at the published
precision), \DFD\ predicts a one-way asymmetry:
\begin{align}
\Delta\psi_{\rm GP-A} &= \frac{\rSch^\oplus}{\Rearth}
\left(1 - \frac{\Rearth}{\Rearth + h}\right) \nonumber\\
&= 1.39\times 10^{-9}\times\frac{10^4}{6371 + 10^4\text{ km}} \nonumber\\
&\approx 8.5\times 10^{-10},
\end{align}
giving $\delta t_{\rm NR}^{\rm GP-A} = 2\times 8.5\times 10^{-10}
\times (10^4\text{ km}/c) \approx 0.057$~ns.

\subsection{Hafele--Keating experiment (1972)}

The Hafele--Keating experiment~\cite{HafeleKeating1972a,HafeleKeating1972b}
transported cesium clocks on commercial aircraft.  Results:
\begin{itemize}
\item Eastbound: $-59\pm 10$~ns (GR: $-40\pm 23$~ns)
\item Westbound: $+273\pm 7$~ns (GR: $+275\pm 21$~ns)
\end{itemize}

\textbf{DFD prediction:}  The \DFD\ nonreciprocal effect for
aircraft at $h \approx 10$~km:
\begin{equation}
\Delta\psi_{\rm HK} = \frac{g\,h}{c^2}\times 2
= \frac{9.8\times 10^4}{9\times 10^{16}} \approx 2.2\times 10^{-12}.
\end{equation}

For a circumnavigation at $\sim 300$~m/s over $\sim 50$~hours,
the accumulated nonreciprocal effect is:
\begin{equation}
\delta t_{\rm NR}^{\rm HK} = \frac{2\Delta\psi_{\rm HK}}{c}
\times v_{\rm aircraft}\times T \approx 0.02~\text{ns}.
\end{equation}

This is well below the $\pm 10$~ns measurement uncertainty.
\DFD\ predicts the Hafele--Keating result is indistinguishable
from GR at the achieved precision.

\subsection{One-way speed of light: Krisher et al.\ (1990)}

Krisher, Maleki, Lutes, and collaborators~\cite{Krisher1990}
performed a one-way speed-of-light measurement using two hydrogen
masers separated by 21~km of optical fiber, achieving a limit
$|\Delta c/c| < 3.5\times 10^{-7}$.

\textbf{DFD prediction:}  For a 21~km horizontal fiber with
negligible height difference:
\begin{equation}
\delta t_{\rm NR}^{\rm Krisher} = \frac{2\Delta\psi}{c}\times\frac{L}{c}
\approx 0,
\end{equation}
since $\Delta\psi \approx 0$ for a horizontal link.  \DFD\ is
consistent with the null result.  The \DFD\ signal would require
a vertical component.

\subsection{One-way light speed: fiber-optic tests (Nagel et al.\ 2015)}

Nagel et al.~\cite{Nagel2015} used crossed optical resonators to
test Lorentz invariance, achieving limits on the anisotropy of the
one-way speed of light at the $10^{-18}$ level.

\textbf{DFD prediction:}  These experiments test for
direction-dependent speed at a \emph{single} spatial point.
\DFD\ predicts no anisotropy at a point---the one-way speed
$c_1 = c\,e^{-\psi}$ depends on position but not direction.
\DFD\ is consistent with the null result.

\subsection{Sagnac effect: precision measurements}

The Sagnac effect for a rotating interferometer of area $A$ is
\begin{equation}
\delta t_{\rm Sagnac} = \frac{4\boldsymbol{\omega}\cdot\mathbf{A}}{c^2},
\end{equation}
which is a topological effect (depends on enclosed area) and is
identical in GR and \DFD.

\textbf{DFD prediction:}  The Sagnac effect is reproduced exactly.
The \DFD\ nonreciprocal effect is \emph{distinct}---it depends on
$\Delta\psi$ (height difference), not enclosed area.  The two
effects are geometrically orthogonal and experimentally separable.

Ring-laser gyroscopes (e.g., the G-ring at Wettzell~\cite{Schreiber2004})
measure the Sagnac effect to $10^{-9}\,\omega_\oplus$ precision.
Since these are horizontal instruments ($\Delta h \approx 0$), \DFD\
predicts zero additional effect, consistent with observations.

\subsection{Fiber-optic time transfer: REFIMEVE+ (Lisdat et al.\ 2016)}

The REFIMEVE+ network~\cite{Lisdat2016} connects European optical
clocks via phase-stabilised fibers, achieving $10^{-18}$-level
comparisons.  The PTB--LNE-SYRTE link (1400~km, $\Delta h \approx 100$~m)
showed gravitational redshift consistent with GR.

\textbf{DFD prediction:}  The standard bidirectional technique
cancels the nonreciprocal term by construction (round-trip
phase stabilisation).  The gravitational redshift measurement
is the symmetric component, which \DFD\ reproduces exactly.

The nonreciprocal component would appear as a persistent
asymmetry between independently measured one-way delays:
\begin{align}
\delta t_{\rm NR}^{\rm PTB-SYRTE} &= \frac{2g\,\Delta h}{c^3}
\times L_{\rm fiber} \nonumber\\
&= \frac{2\times 9.8\times 100}{2.7\times 10^{25}}
\times 1.4\times 10^9 \nonumber\\
&= 1.02\times 10^{-16}~\text{s} = 0.10~\text{fs}.
\label{eq:PTB_SYRTE}
\end{align}

This is at the edge of current sensitivity.

\subsection{Delva et al.\ (2017): Fiber test of special relativity}

Delva et al.~\cite{Delva2017} used the LNE-SYRTE fiber network
to test the isotropy of the speed of light, finding consistency
with special relativity at the $10^{-15}$ level.

\textbf{DFD prediction:}  The experiment tested \emph{isotropy}
(direction-dependence at a single point), not nonreciprocity
(direction-of-propagation dependence between two points).
\DFD\ predicts no isotropy violation, consistent with the result.

\subsection{VLBI timing residuals (Titov et al.\ 2011)}

Titov, Lambert, and Gontier~\cite{Titov2011} measured the secular
aberration drift using 30 years of VLBI data and found a dipole
signal consistent with the galactocentric acceleration of the
Solar System: $a_{\rm gal} = (5.2\pm 0.2)\times 10^{-10}$~m/s$^2$
toward the Galactic centre.

\textbf{DFD prediction:}  VLBI uses differenced delays between
two ground stations observing the same source.  For a baseline
$\mathbf{B}$ with height difference $\Delta h_B$:
\begin{equation}
\delta t_{\rm NR}^{\rm VLBI} \approx \frac{2g\,\Delta h_B}{c^3}
\times \frac{L_{\rm baseline}}{c}\times c = \frac{2g\,\Delta h_B}{c^2}.
\end{equation}

For a transcontinental baseline with $\Delta h_B = 1000$~m:
$\delta t_{\rm NR}^{\rm VLBI} \approx 2.2\times 10^{-14}$~s $= 22$~fs.
This is well below VLBI timing precision ($\sim 5$~ps), so
\DFD\ predicts no detectable anomaly in current VLBI data.

\subsection{Pulsar timing arrays (NANOGrav, PPTA, EPTA)}

Pulsar timing arrays~\cite{NANOGrav2023} achieve $\sim 100$~ns
residuals for millisecond pulsars.  These are inherently one-way
measurements (pulsar to Earth).

\textbf{DFD prediction:}  The Galactic $\psi$-gradient between
pulsar and Earth introduces a direction-dependent one-way correction:
\begin{equation}
\delta t_{\rm NR}^{\rm PTA} = \frac{2}{c^2}\int_{\rm pulsar}^{\rm Earth}
\nabla\psi_{\rm MW}\cdot d\boldsymbol{\ell}\times\frac{D_{\rm psr}}{c},
\end{equation}
where $D_{\rm psr}$ is the pulsar distance.  For a pulsar at
$D = 1$~kpc in the Galactic plane with the full Milky Way potential:
\begin{equation}
\delta t_{\rm NR}^{\rm PTA} \sim \frac{2\Delta\psi_{\rm MW}}{c}
\times \frac{D}{c} \sim \frac{2\times 10^{-6}}{c}\times 10^{11}
\sim 0.7~\text{ms}.
\end{equation}

This is a \emph{static} offset that would be absorbed into the
pulsar distance/dispersion measure fit.  The \emph{time-varying}
component (from Earth's orbital motion in the Galactic potential)
has period $\sim 250$~Myr and is undetectable.

However, the annual modulation of the Galactic $\psi$-gradient
projection (from Earth's orbital motion around the Sun changing
the line-of-sight to the pulsar) produces:
\begin{equation}
\delta t_{\rm NR}^{\rm PTA,annual} \sim |\nabla\psi_{\rm MW}|
\times 2\,\text{AU}/c \times D_{\rm psr}/c
\sim 10^{-18}\times 10^3\times 10^{11}/c \sim 30~\text{ps}.
\end{equation}

This is below current PTA sensitivity but could emerge with
next-generation timing precision.

\subsection{Lunar laser ranging (Williams et al.\ 2004)}

Lunar laser ranging (LLR)~\cite{Williams2004} is a \emph{two-way}
measurement: laser pulse goes to the Moon and returns.  Achieved
precision: $\sim 1$~cm ($\sim 33$~ps in round-trip time).

\textbf{DFD prediction:}  By Corollary~1, the nonreciprocal term
cancels in the round-trip measurement.  \DFD\ predicts zero
anomaly beyond GR in LLR data.

The one-way nonreciprocal correction would be:
\begin{align}
\Delta\psi_{\rm Moon} &= \frac{\rSch^\oplus}{\Rearth}
\left(1 - \frac{\Rearth}{R_{\rm Moon}}\right) \nonumber\\
&= 1.39\times 10^{-9}\times\left(1 - \frac{6371}{384400}\right)
= 1.37\times 10^{-9},
\end{align}
giving $\delta t_{\rm NR}^{\rm Moon} = 2\times 1.37\times 10^{-9}
\times 1.28\text{ s} = 3.5$~ns one-way.  But this is invisible
in the two-way measurement.

\subsection{Shapiro delay: Cassini (Bertotti et al.\ 2003)}

The Cassini solar conjunction experiment~\cite{Bertotti2003}
measured the Shapiro delay to test $\gamma = 1 + (2.1\pm 2.3)
\times 10^{-5}$.  This used two-way ranging with Doppler tracking.

\textbf{DFD prediction:}  Two-way measurement---nonreciprocal
term cancels.  \DFD\ reproduces $\gamma = 1$ exactly and is
consistent with the Cassini result.

\subsection{Planetary radar ranging: Mercury, Venus}

Planetary radar ranging is two-way (radar echo).  The $\gamma$
measurements from Mercury and Venus inferior conjunctions
constrain $|\gamma - 1| < 3\times 10^{-4}$~\cite{Shapiro1971}.

\textbf{DFD prediction:}  Two-way---nonreciprocal term cancels.
Consistent with all radar ranging results.

\subsection{Atomic interferometer tests (M\"{u}ller et al.\ 2010)}

M\"{u}ller, Peters, and Chu~\cite{Mueller2010} measured the
gravitational redshift using atom interferometry, claiming
$7\times 10^{-9}$ precision.

\textbf{DFD prediction:}  Atom interferometry measures the phase
difference accumulated by a matter wave traversing different
gravitational potentials.  This is a local measurement at a
single laboratory, not a one-way propagation measurement between
separated clocks.  \DFD\ reproduces the standard gravitational
phase shift exactly.

\subsection{Summary of literature predictions}

\begin{table*}[t]
\caption{Complete \DFD\ predictions for 15 precision timing experiments.
``NR'' = nonreciprocal correction; ``2-way'' = cancels in round-trip.
All predictions are consistent with published results.}
\label{tab:literature}
\begin{ruledtabular}
\begin{tabular}{llccc}
Experiment & Type & DFD NR signal & Measurable? & Consistent? \\
\hline
GPS residuals (Senior 2008) & One-way & $0.14$--$0.21$ ns & Yes & Yes \\
GPS time transfer (Petit \& Jiang 2008) & One-way & $\sim 0.05$ ns jumps & Marginal & Yes \\
Gravity Probe A (Vessot 1980) & Two-way & Cancels & N/A & Yes \\
Hafele--Keating (1972) & Transported clock & $0.02$ ns & Below noise & Yes \\
Krisher et al.\ (1990) & Horizontal fiber & $\approx 0$ & N/A & Yes \\
Nagel et al.\ (2015) & Local resonator & $0$ & N/A & Yes \\
Sagnac (Schreiber 2004) & Ring interferometer & $0$ (horizontal) & N/A & Yes \\
REFIMEVE+ (Lisdat 2016) & Bidirectional fiber & $0.10$ fs & At edge & Yes \\
Delva et al.\ (2017) & Isotropy test & $0$ & N/A & Yes \\
VLBI (Titov 2011) & Baseline differencing & $22$ fs & Below noise & Yes \\
Pulsar timing arrays & One-way (static) & Absorbed in fit & N/A & Yes \\
Lunar laser ranging & Two-way & Cancels & N/A & Yes \\
Cassini (Bertotti 2003) & Two-way & Cancels & N/A & Yes \\
Planetary radar & Two-way & Cancels & N/A & Yes \\
Atom interferometry & Local phase & $0$ & N/A & Yes \\
\end{tabular}
\end{ruledtabular}
\end{table*}

%=============================================================================
\section{The Hubble Tension: Rigorous Derivation}\label{sec:hubble}
%=============================================================================

This section derives the connection between the local $\psi$-gradient
and the Hubble tension with full rigor, starting from a realistic
Milky Way mass model.

\subsection{Milky Way mass model}

We adopt the standard three-component Milky Way model
(McMillan 2017~\cite{McMillan2017}):

\textbf{1. Bulge} (Hernquist profile):
\begin{equation}
\Phi_{\rm bulge}(r) = -\frac{GM_{\rm b}}{r + a_{\rm b}},
\quad M_{\rm b} = 1.0\times 10^{10}\,\Msun,\;\; a_{\rm b} = 0.6~\text{kpc}.
\end{equation}

\textbf{2. Disk} (Miyamoto--Nagai):
\begin{equation}
\Phi_{\rm disk}(R,z) = -\frac{GM_{\rm d}}
{\sqrt{R^2 + (a_{\rm d} + \sqrt{z^2 + b_{\rm d}^2})^2}},
\end{equation}
with $M_{\rm d} = 5.0\times 10^{10}\,\Msun$,
$a_{\rm d} = 3.0$~kpc, $b_{\rm d} = 0.28$~kpc.

\textbf{3. Dark matter halo} (NFW profile):
\begin{equation}
\Phi_{\rm NFW}(r) = -\frac{GM_{\rm vir}}{r\,[\ln(1+c)-c/(1+c)]}
\ln\!\left(1 + \frac{r}{r_s}\right),
\end{equation}
with virial mass $M_{\rm vir} = 1.3\times 10^{12}\,\Msun$,
concentration $c = 12$, scale radius $r_s = r_{\rm vir}/c = 20$~kpc.

\textbf{Note:} In \DFD, the ``dark matter halo'' is not physical
dark matter but rather the nonlinear $\mu$-function regime of the
$\psi$-field.  For this calculation, we use the effective potential
that reproduces the observed rotation curve, regardless of its
microscopic origin.

\subsection{The $\psi$-field at the Solar position}

The Sun is at $R_\odot = 8.2$~kpc from the Galactic centre in
the disk plane ($z \approx 25$~pc $\approx 0$).  The total
Newtonian potential:
\begin{align}
\Phi_\odot &= \Phi_{\rm bulge}(8.2) + \Phi_{\rm disk}(8.2,0)
+ \Phi_{\rm NFW}(8.2) \nonumber\\
&\approx (-1.06 - 2.14 - 2.99)\times 10^{10}~\text{m}^2/\text{s}^2
\nonumber\\
&= -6.19\times 10^{10}~\text{m}^2/\text{s}^2.
\end{align}

Converting to $\psi$:
\begin{equation}
\psi_\odot = \frac{2\Phi_\odot}{c^2}
= \frac{2\times(-6.19\times 10^{10})}{9\times 10^{16}}
= -1.376\times 10^{-6}.
\label{eq:psi_solar}
\end{equation}

\subsection{The $\psi$-gradient at the Solar position}

The radial gradient (pointing toward the Galactic centre):
\begin{equation}
\frac{\partial\psi}{\partial R}\bigg|_\odot
= \frac{2}{c^2}\frac{\partial\Phi}{\partial R}\bigg|_\odot
= -\frac{2\,a_{\rm gal}}{c^2},
\end{equation}
where $a_{\rm gal} = v_c^2/R_\odot$ is the centripetal acceleration.

Using $v_c = 232$~km/s and $R_\odot = 8.2$~kpc:
\begin{equation}
a_{\rm gal} = \frac{(2.32\times 10^5)^2}{2.53\times 10^{20}}
= 2.13\times 10^{-10}~\text{m/s}^2.
\end{equation}

The $\psi$-gradient magnitude:
\begin{equation}
\left|\nabla\psi\right|_\odot = \frac{2\,a_{\rm gal}}{c^2}
= \frac{2\times 2.13\times 10^{-10}}{9\times 10^{16}}
= 4.73\times 10^{-27}~\text{m}^{-1}.
\label{eq:grad_psi_MW}
\end{equation}

\subsection{Laniakea supercluster contribution}

The Solar System is falling toward the Great Attractor region
at $v_{\rm pec} \approx 630$~km/s.  The Laniakea supercluster
has a characteristic mass $M_{\rm Lan} \sim 10^{17}\,\Msun$
at distance $D_{\rm Lan} \sim 80$~Mpc.

The supercluster $\psi$-gradient:
\begin{equation}
|\nabla\psi_{\rm Lan}| = \frac{2GM_{\rm Lan}}{c^2 D_{\rm Lan}^2}
= \frac{2\times 6.674\times 10^{-11}\times 2\times 10^{47}}
{9\times 10^{16}\times(2.47\times 10^{24})^2}.
\end{equation}

Computing:
\begin{equation}
|\nabla\psi_{\rm Lan}| = \frac{2.67\times 10^{37}}{5.49\times 10^{65}}
= 4.86\times 10^{-29}~\text{m}^{-1}.
\end{equation}

This is subdominant to the Milky Way contribution by a factor
$\sim 100$.

\subsection{The $\psi$-screen and luminosity distance modification}

In \DFD, a photon propagating from redshift $z$ to the observer
accumulates a $\psi$-screen:
\begin{equation}
\Delta\psi(z,\hat{n}) = \int_0^{\chi(z)} \nabla\psi\cdot\hat{n}\,d\chi'
+ \psi_{\rm cosmo}(z),
\end{equation}
where the first term is the directional (anisotropic) component
and $\psi_{\rm cosmo}(z)$ is the isotropic cosmological component.

The modified luminosity distance:
\begin{equation}
D_L^{\rm DFD}(z,\hat{n}) = D_L^{\rm true}(z)\,
\exp\!\left[\Delta\psi(z,\hat{n})\right].
\end{equation}

The inferred Hubble constant from a local distance measurement
(using standard candles at distance $d$):
\begin{equation}
H_0^{\rm inferred}(\hat{n}) = H_0^{\rm true}\,
\exp\!\left[-\Delta\psi_{\rm local}(\hat{n})\right].
\end{equation}

\subsection{Computing the Hubble shift: monopole component}

The monopole $\psi$-screen is the \emph{potential difference}
between our location and a ``fair sample'' of the cosmos.
A galaxy sitting at the bottom of its potential well has
$\psi_{\rm gal} < 0$, while the cosmic mean has $\psi = 0$.

The Solar System sits in the combined potential well of:
\begin{enumerate}
\item Milky Way: $\psi_{\rm MW} = -1.38\times 10^{-6}$
\item Local Group: $\psi_{\rm LG} \approx -4\times 10^{-7}$
      (from M31 and group potential)
\item Virgo Supercluster: $\psi_{\rm Virgo} \approx -3\times 10^{-7}$
\item Laniakea: $\psi_{\rm Lan} \approx -2\times 10^{-7}$
\end{enumerate}

Total local $\psi$-well:
\begin{equation}
\psi_{\rm local} \approx -2.3\times 10^{-6}.
\label{eq:psi_local}
\end{equation}

The shift in inferred $H_0$:
\begin{equation}
\frac{\delta H_0}{H_0}\bigg|_{\rm monopole}
= -\psi_{\rm local}
\approx +2.3\times 10^{-6}.
\label{eq:H0_monopole_psi}
\end{equation}

This monopole shift is only $\sim 2\times 10^{-4}\%$ of $H_0$,
far too small to explain the $\sim 8\%$ Hubble tension directly
through the potential well alone.

\subsection{The accumulated $\psi$-screen along the line of sight}

The crucial insight is that the $\psi$-screen accumulates
\emph{along the entire light path}, not just at the endpoints.
For a photon traversing a distance $\chi$ through the cosmic
density field, the effective $\psi$-screen is:
\begin{equation}
\Delta\psi_{\rm eff}(z) = \frac{2}{c^2}\int_0^{\chi(z)}
\delta\Phi(\mathbf{x}(\chi'),\chi')\,d\chi',
\end{equation}
where $\delta\Phi$ is the gravitational potential fluctuation
relative to the cosmic mean.

This is closely related to the Integrated Sachs--Wolfe (ISW)
effect and gravitational lensing convergence.  The key difference
in \DFD\ is that $\psi$ affects the \emph{one-way} travel time
(not just the photon energy), producing a distance bias rather
than just a temperature shift.

For standard candles at $z \sim 0.05$ (the SH0ES Cepheid sample),
the line-of-sight $\psi$-screen depends on the cosmic web
structure.  The variance of the $\psi$-screen is:
\begin{equation}
\sigma_\psi^2(z) = \frac{4}{c^4}\int \frac{dk}{k}\,
\frac{k^3 P_\Phi(k)}{2\pi^2}\,W^2(k,\chi(z)),
\end{equation}
where $P_\Phi(k)$ is the gravitational potential power spectrum
and $W$ is a window function.

Using the standard $\Lambda$CDM power spectrum at $z \sim 0.05$:
\begin{equation}
\sigma_\psi(z = 0.05) \sim 3\times 10^{-5}.
\end{equation}

\subsection{Computing the apparent $H_0$}

The critical question: can the $\psi$-screen produce a
\emph{systematic} bias in $H_0$ of $\sim 8\%$?

The SH0ES distance ladder measures $H_0$ using Cepheids in
galaxies at $d \lesssim 40$~Mpc.  These host galaxies sit in
their own potential wells.  If the \DFD\ $\psi$-screen
systematically biases the Cepheid period-luminosity relation
or the SNe~Ia standardisation, the effect is:
\begin{equation}
H_0^{\rm SH0ES} = H_0^{\rm true}\,
\exp\!\left[-\langle\Delta\psi\rangle_{\rm hosts}\right].
\end{equation}

For the host galaxy potential wells, $\psi_{\rm host} \sim -10^{-6}$
to $-10^{-5}$, yielding corrections of order $10^{-6}$ to $10^{-5}$.

\textbf{Honest assessment:} The simple $\psi$-well mechanism gives
$\delta H_0/H_0 \sim 10^{-6}$ to $10^{-5}$, which is
\emph{far too small} to explain the $\sim 8\%$ Hubble tension
through the Newtonian-limit $\psi$-field alone.

\subsection{The MOND-regime enhancement}

The previous calculation used the weak-field (Newtonian) relation
$\psi = 2\Phi/c^2$.  In \DFD, the nonlinear $\mu$-function
modifies this in the low-acceleration regime.  The \DFD\ field
equation:
\begin{equation}
\nabla\cdot[\mu(|\nabla\psi|/\astar)\nabla\psi]
= \frac{8\pi G}{c^2}\rho,
\end{equation}
with $\mu(x) = x/(1+x)$ (the derived simple form), gives an
\emph{enhanced} $\psi$ in the deep-field regime where
$|\nabla\psi|/\astar \ll 1$.

In the deep-field limit, $\mu(x) \approx x$, so:
\begin{equation}
|\nabla\psi|^2 \approx \astar\times\frac{8\pi G}{c^2}\rho
\times R_{\rm eff},
\end{equation}
which gives $|\nabla\psi| \sim \sqrt{\astar\times g_N}$
(the MOND acceleration).

For the Milky Way at large radii ($r \gtrsim 50$~kpc), the
$\psi$-field falls off as $\psi \sim -\sqrt{GM\,\azero}\,\ln(r/r_0)/c^2$
rather than $\psi \sim -GM/(c^2 r)$.  The logarithmic fall-off
means $\psi$ is much larger (more negative) at large radii than
the Newtonian prediction.

Integrating the MOND-regime $\psi$-field from the Milky Way's
virial radius ($r_{\rm vir} \approx 280$~kpc) to
$r_{\rm eff} \sim 1$~Mpc:
\begin{equation}
\psi_{\rm MOND}(r_{\rm eff}) \approx
-\frac{\sqrt{GM_{\rm MW}\,\azero}}{c^2}\ln\!\left(\frac{r_{\rm eff}}{r_0}\right).
\end{equation}

With $M_{\rm MW} = 1.3\times 10^{12}\,\Msun$ and
$\azero = 1.2\times 10^{-10}$~m/s$^2$:
\begin{align}
\sqrt{GM_{\rm MW}\,\azero} &= \sqrt{6.674\times 10^{-11}\times 2.59\times 10^{42}\times 1.2\times 10^{-10}} \nonumber\\
&= \sqrt{2.07\times 10^{22}} = 4.55\times 10^{11}~\text{m}^2/\text{s}^2.
\end{align}

Wait---this should be $\sqrt{GM\azero}$ which has units of
m$^2$/s$^2$ divided by length... Let us be more careful.

The deep-MOND acceleration at radius $r$ is:
\begin{equation}
a_{\rm MOND}(r) = \sqrt{a_N(r)\times\azero}
= \sqrt{\frac{GM\azero}{r^2}} = \frac{\sqrt{GM\azero}}{r}.
\end{equation}

The potential in the deep-MOND regime:
\begin{equation}
\Phi_{\rm MOND}(r) = -\int_r^\infty a_{\rm MOND}(r')\,dr'
= -\sqrt{GM\azero}\ln\!\left(\frac{r_\infty}{r}\right).
\end{equation}

With $\sqrt{GM_{\rm MW}\azero} = \sqrt{8.66\times 10^{31}\times 1.2\times 10^{-10}} = \sqrt{1.04\times 10^{22}} = 3.22\times 10^{11}$~(m/s)$^2$... but this has wrong units.

Let us compute properly:
\begin{align}
GM_{\rm MW} &= 6.674\times 10^{-11}\times 1.3\times 10^{12}\times 1.989\times 10^{30} \nonumber\\
&= 6.674\times 10^{-11}\times 2.586\times 10^{42} \nonumber\\
&= 1.726\times 10^{32}~\text{m}^3/\text{s}^2. \\
\sqrt{GM_{\rm MW}\,\azero} &= \sqrt{1.726\times 10^{32}\times 1.2\times 10^{-10}} \nonumber\\
&= \sqrt{2.07\times 10^{22}} = 4.55\times 10^{11}~\text{m}^2/\text{s}^2.
\end{align}

So indeed $[\sqrt{GM\azero}] = \text{m}^2/\text{s}^2$ and
$\Phi_{\rm MOND} = -\sqrt{GM\azero}\ln(r_\infty/r)$.  Then:
\begin{equation}
\psi_{\rm MOND}(r) = \frac{2\Phi_{\rm MOND}}{c^2}
= -\frac{2\sqrt{GM\azero}}{c^2}\ln\!\left(\frac{r_\infty}{r}\right).
\end{equation}

The ``depth'' of the MOND $\psi$-well at the Solar position
(with $r_\infty$ set by the surrounding cosmic density, roughly
$r_\infty \sim 1$--$3$~Mpc, and $R_\odot = 8.2$~kpc):
\begin{align}
|\psi_{\rm MOND}| &= \frac{2\times 4.55\times 10^{11}}{9\times 10^{16}}
\times\ln\!\left(\frac{3\times 10^6}{8.2}\right) \nonumber\\
&= 1.01\times 10^{-5}\times\ln(3.66\times 10^5) \nonumber\\
&= 1.01\times 10^{-5}\times 12.81 \nonumber\\
&= 1.29\times 10^{-4}.
\label{eq:psi_MOND_well}
\end{align}

\textbf{This is 100 times larger than the Newtonian estimate.}

\subsection{The apparent $H_0$ with MOND-regime $\psi$-screen}

Using the MOND-regime $\psi$-well:
\begin{equation}
\frac{\delta H_0}{H_0}\bigg|_{\rm DFD}
= -\psi_{\rm local}^{(\rm MOND)} \approx +1.3\times 10^{-4},
\end{equation}
or $\delta H_0 \approx 0.009$~km/s/Mpc.  Still too small by a
factor $\sim 600$.

\subsection{Full cosmological $\psi$-screen: the distance-duality argument}

The monopole Hubble shift cannot come from the local potential well
alone.  However, in \DFD\, there is a more subtle effect: the
\emph{cosmic expansion} itself is reinterpreted as a change in the
cosmic $\psi$-field.

In \DFD\ cosmology, the Friedmann equation becomes:
\begin{equation}
H^2 = \frac{8\pi G}{3}\rho_{\rm total}
- \frac{\kappa}{c^2}\left(\frac{\dot\psi_{\rm cosmo}}{H}\right)^2,
\end{equation}
where the second term is the $\psi$-field contribution to the
effective expansion rate.  The apparent acceleration ($q < 0$)
arises not from a cosmological constant but from the $\psi$-field's
nonlinear self-interaction in the deep-field (MOND) cosmological
regime.

The full derivation of the cosmological $\psi$-screen and its
effect on $H_0$ requires the complete DFD cosmological solution,
which is beyond the scope of this GPS-focused paper.  We note
the following bounds:

\textbf{Lower bound (local potential well only):}
$\delta H_0/H_0 \sim 10^{-4}$ (Eq.~\ref{eq:psi_MOND_well}).

\textbf{Upper bound (if DFD replaces $\Lambda$ entirely):}
The dark energy contribution to $H_0$ in $\Lambda$CDM is
$H_0^2 = H_\infty^2 + (8\pi G/3)\rho_{m,0}$, where
$H_\infty = H_0\sqrt{\Omega_\Lambda} \approx 56$~km/s/Mpc.
If the \DFD\ $\psi$-screen modifies the distance--redshift
relation in a way that mimics a different $\Omega_\Lambda$,
the apparent $H_0$ can shift by the full tension:
$\delta H_0 \sim 5.6$~km/s/Mpc.

\textbf{Best estimate (dipole component):}
The directional Hubble dipole observed by Secrest et al.~\cite{Secrest2021}
has amplitude $\sim 2\%$, corresponding to $\delta H_0 \sim 1.4$~km/s/Mpc.
The \DFD\ $\psi$-gradient at the Solar position (Eq.~\ref{eq:grad_psi_MW})
integrated over the distance to the SH0ES calibrator galaxies
($d \sim 20$--$40$~Mpc):
\begin{align}
\Delta\psi_{\rm dipole} &= |\nabla\psi|_\odot \times d \nonumber\\
&= 4.73\times 10^{-27}\times 10^{24} \approx 5\times 10^{-3}.
\label{eq:dipole_psi}
\end{align}

Wait---this uses the Newtonian $|\nabla\psi|$.  In the MOND regime,
the gradient at the Solar position in the Galactic potential is:
\begin{equation}
|\nabla\psi|_\odot = \frac{2\,a_{\rm gal}}{c^2}
= \frac{2\times 2.13\times 10^{-10}}{9\times 10^{16}}
= 4.73\times 10^{-27}~\text{m}^{-1}.
\end{equation}

But outside the Milky Way, the gradient falls off.  The relevant
integral is the line-of-sight integral through the Galactic
potential:
\begin{align}
\Delta\psi_{\rm LOS} &= \int_0^d |\nabla\psi_{\rm MW}|(r,\hat{n})\,dr \nonumber\\
&\approx \psi_\odot - \psi(\text{intergalactic})
\approx 1.3\times 10^{-4}~\text{(MOND)}.
\end{align}

This gives:
\begin{equation}
\frac{\delta H_0}{H_0}\bigg|_{\rm dipole,max}
\sim 1.3\times 10^{-4}\times\cos\theta_{\rm GC},
\end{equation}
where $\theta_{\rm GC}$ is the angle from the Galactic centre
direction.  The dipole amplitude is $\sim 0.013\%$, far below
the observed $\sim 2\%$.

\textbf{Conclusion:} The simple \DFD\ Newtonian/MOND $\psi$-well
around the Milky Way cannot explain the full Hubble tension.
The mechanism would require the full cosmological $\psi$-screen
calculation, which modifies the distance--redshift relation
globally.  The local $\psi$-gradient produces a dipole of order
$0.01\%$, which is consistent with current dipole measurements
but does not resolve the monopole tension.

We report this honestly as a negative result for the simplest
version of the Hubble tension connection.  The full cosmological
calculation remains an open problem in DFD.

%=============================================================================
\section{Deep-Space Ranging Corrections}\label{sec:deep_space}
%=============================================================================

We compute the \DFD\ one-way nonreciprocal correction for each
major deep-space mission.  The dominant $\psi$-gradient is the
\emph{Sun's} field, since inter-planetary links traverse a
significant fraction of the Solar potential.

\subsection{General formula}

For two bodies at heliocentric distances $r_1$ and $r_2$, the
solar $\Delta\psi$:
\begin{equation}
\Delta\psi_\odot = \frac{\rSch^\odot}{r_1} - \frac{\rSch^\odot}{r_2}
= \rSch^\odot\left(\frac{1}{r_1} - \frac{1}{r_2}\right),
\end{equation}
where $\rSch^\odot = 2G\Msun/c^2 = 2.954$~km.

The one-way light travel time is $T_{\rm 1-way} = d_{12}/c$,
where $d_{12}$ is the inter-body distance.

The nonreciprocal correction:
\begin{equation}
\delta t_{\rm NR}^{(\odot)} = \frac{2\Delta\psi_\odot}{c}
\times T_{\rm 1-way}.
\label{eq:NR_solar_ds}
\end{equation}

In addition, each body has its own local $\psi$ from its own mass:
\begin{equation}
\Delta\psi_{\rm local} = \frac{\rSch^{(\rm body)}}{r_{\rm surface/orbit}}
- \frac{\rSch^\oplus}{\Rearth}.
\end{equation}

\subsection{Mars--Earth (all missions)}

Mars orbital parameters: $a_{\rm Mars} = 1.524$~AU $= 2.279\times 10^{11}$~m.

Solar $\Delta\psi$ (Earth to Mars):
\begin{align}
\Delta\psi_\odot^{(\rm Mars)} &= 2.954\times 10^3
\left(\frac{1}{1.496\times 10^{11}} - \frac{1}{2.279\times 10^{11}}\right) \nonumber\\
&= 2.954\times 10^3\times(6.685 - 4.388)\times 10^{-12} \nonumber\\
&= 2.954\times 10^3\times 2.297\times 10^{-12} \nonumber\\
&= 6.785\times 10^{-9}.
\end{align}

\textbf{At closest approach} ($d \approx 0.52$~AU $= 7.78\times 10^{10}$~m,
$T_{\rm 1-way} = 259$~s):
\begin{equation}
\delta t_{\rm NR}^{\rm Mars,close} = \frac{2\times 6.785\times 10^{-9}}{3\times 10^8}
\times 259 = 11.7~\text{ns}.
\end{equation}

\textbf{At opposition} ($d \approx 0.52$~AU): Same as above.

\textbf{At conjunction} ($d \approx 2.52$~AU $= 3.77\times 10^{11}$~m,
$T_{\rm 1-way} = 1258$~s):
\begin{equation}
\delta t_{\rm NR}^{\rm Mars,conj} = \frac{2\times 6.785\times 10^{-9}}{3\times 10^8}
\times 1258 = 56.9~\text{ns}.
\end{equation}

\textbf{Mars missions and predictions:}
\begin{itemize}
\item \textbf{Viking (1976):} Used two-way Doppler ranging.
      Nonreciprocal term cancels.
\item \textbf{Mars Pathfinder (1997):} Two-way Doppler.  Cancels.
\item \textbf{MRO (2006--present):} One-way and two-way modes.
      One-way mode prediction: $\delta t_{\rm NR} = 12$--$57$~ns
      depending on orbital geometry.
\item \textbf{InSight (2018--2022):} RISE experiment used
      two-way X-band Doppler.  Cancels.
\item \textbf{Perseverance (2021--present):} Two-way.  Cancels.
\end{itemize}

\subsection{Jupiter--Earth (Juno)}

Jupiter: $a_{\rm Jup} = 5.203$~AU $= 7.783\times 10^{11}$~m.

Solar $\Delta\psi$:
\begin{align}
\Delta\psi_\odot^{(\rm Jup)} &= 2.954\times 10^3
\times(6.685 - 1.285)\times 10^{-12} \nonumber\\
&= 2.954\times 10^3\times 5.400\times 10^{-12}
= 1.595\times 10^{-8}.
\end{align}

At closest approach ($d \approx 4.2$~AU, $T_{\rm 1-way} = 2097$~s):
\begin{equation}
\delta t_{\rm NR}^{\rm Jup,close} = \frac{2\times 1.595\times 10^{-8}}{3\times 10^8}
\times 2097 = 223~\text{ns}.
\end{equation}

At conjunction ($d \approx 6.2$~AU, $T_{\rm 1-way} = 3094$~s):
\begin{equation}
\delta t_{\rm NR}^{\rm Jup,conj} = \frac{2\times 1.595\times 10^{-8}}{3\times 10^8}
\times 3094 = 329~\text{ns}.
\end{equation}

Juno uses two-way coherent Doppler (X-band uplink, X/Ka-band
downlink), so the nonreciprocal term cancels in the primary
science data.  However, the one-way downlink telemetry timing
carries the \DFD\ prediction of $\sim 220$--$330$~ns nonreciprocal
offset.

\subsection{Saturn--Earth (Cassini)}

Saturn: $a_{\rm Sat} = 9.537$~AU $= 1.427\times 10^{12}$~m.

Solar $\Delta\psi$:
\begin{align}
\Delta\psi_\odot^{(\rm Sat)} &= 2.954\times 10^3
\times(6.685 - 0.701)\times 10^{-12} \nonumber\\
&= 2.954\times 10^3\times 5.984\times 10^{-12}
= 1.767\times 10^{-8}.
\end{align}

At closest approach ($d \approx 8.5$~AU, $T_{\rm 1-way} = 4246$~s):
\begin{equation}
\delta t_{\rm NR}^{\rm Sat,close} = \frac{2\times 1.767\times 10^{-8}}{3\times 10^8}
\times 4246 = 500~\text{ns}.
\end{equation}

Cassini used two-way coherent ranging for the $\gamma$ measurement
(Bertotti et al.\ 2003), so the nonreciprocal term cancels in that
result.

\subsection{Pluto--Earth (New Horizons)}

Pluto: $a_{\rm Pluto} \approx 39.5$~AU $= 5.91\times 10^{12}$~m
(at encounter in 2015).

Solar $\Delta\psi$:
\begin{align}
\Delta\psi_\odot^{(\rm Pluto)} &= 2.954\times 10^3
\times(6.685 - 0.169)\times 10^{-12} \nonumber\\
&= 2.954\times 10^3\times 6.516\times 10^{-12}
= 1.924\times 10^{-8}.
\end{align}

At encounter ($d \approx 32.9$~AU, $T_{\rm 1-way} = 16\,430$~s
$\approx 4.56$~hours):
\begin{equation}
\delta t_{\rm NR}^{\rm Pluto} = \frac{2\times 1.924\times 10^{-8}}{3\times 10^8}
\times 16430 = 2105~\text{ns} \approx 2.1~\mu\text{s}.
\end{equation}

New Horizons uses one-way downlink (no uplink capability during
flyby), making this the most directly testable deep-space
prediction.  The $2.1~\mu$s nonreciprocal correction corresponds
to a range error of $\sim 630$~m.

\subsection{Summary of deep-space predictions}

\begin{table}[t]
\caption{DFD one-way nonreciprocal corrections for deep-space
missions.  All two-way ranging is immune.  ``Close'' = closest
approach; ``Far'' = conjunction/encounter.}
\label{tab:deep_space}
\begin{ruledtabular}
\begin{tabular}{lcccc}
Target & $\Delta\psi_\odot$ & $T_{\rm 1-way}$ & $\delta t_{\rm NR}$ & Ranging mode \\
 & ($10^{-9}$) & (s) & (ns) & \\
\hline
Mars (close) & $6.79$ & 259 & 11.7 & 2-way \\
Mars (conj.) & $6.79$ & 1258 & 56.9 & 2-way \\
Jupiter (close) & $15.95$ & 2097 & 223 & 2-way \\
Jupiter (conj.) & $15.95$ & 3094 & 329 & 2-way \\
Saturn (close) & $17.67$ & 4246 & 500 & 2-way \\
Pluto (enc.) & $19.24$ & 16430 & 2105 & 1-way \\
\end{tabular}
\end{ruledtabular}
\end{table}

\textbf{Key finding:} All precision deep-space ranging experiments
(Viking, Cassini, etc.) used two-way coherent Doppler, rendering
them blind to the nonreciprocal correction.  The \DFD\ prediction
is consistent with all published results.  New Horizons one-way
telemetry timing is the most promising dataset for detecting the
effect.

\subsection{Pioneer anomaly connection}

The Pioneer anomaly~\cite{Anderson2002}---an apparent sunward
acceleration of $(8.74\pm 1.33)\times 10^{-10}$~m/s$^2$---was
eventually attributed to anisotropic thermal radiation
pressure~\cite{Turyshev2012}.  However, \DFD\ would predict a
\emph{one-way timing} anomaly (not a dynamical acceleration)
that could mimic an apparent acceleration in Doppler tracking.

For Pioneer 10 at $r \approx 70$~AU:
\begin{equation}
\delta t_{\rm NR}^{\rm Pioneer} \approx 2.0\times 10^{-8}
\times\frac{70~\text{AU}}{c} \approx 6.8~\mu\text{s}.
\end{equation}

The rate of change (apparent Doppler drift) as the spacecraft
recedes at $v \approx 12$~km/s:
\begin{equation}
\dot{\delta t}_{\rm NR} \approx \frac{2\Delta\psi}{c}\times\frac{v}{c}
\approx 1.3\times 10^{-12}\times 4\times 10^{-5}
\approx 5\times 10^{-17}~\text{s/s}.
\end{equation}

This corresponds to a fractional frequency drift, not an
acceleration.  The apparent acceleration equivalent:
$a_{\rm app} = c\times\ddot{\delta t}_{\rm NR} \sim 10^{-14}$~m/s$^2$,
which is far too small to explain the Pioneer anomaly.  This is
consistent with the thermal explanation being correct.

%=============================================================================
\section{Definitive Experiment: PTB--SYRTE Fiber Link}
\label{sec:experiment}
%=============================================================================

\subsection{Experimental concept}

We propose a definitive test of \DFD\ nonreciprocal timing using
the existing optical fiber link between Physikalisch-Technische
Bundesanstalt (PTB, Braunschweig, Germany) and LNE-SYRTE
(Paris Observatory, France).

\textbf{Key parameters:}
\begin{itemize}
\item Fiber length: $L = 1415$~km
\item Height difference: $\Delta h \approx 25$~m (PTB: 79~m a.s.l.,
      SYRTE: 54~m a.s.l.)\ [Note: actual $\Delta h$ should be
      verified from geodetic surveys]
\item Clock type: Sr optical lattice (PTB) and Sr/Yb (SYRTE)
\item Clock stability: $\sigma_y(\tau) \sim 10^{-16}/\sqrt{\tau/\text{s}}$
      (state of art for optical clocks)
\item Fiber noise suppression: active round-trip phase cancellation,
      residual instability $\sim 10^{-19}$ at $\tau = 10^4$~s
\end{itemize}

\subsection{Predicted signal}

For the PTB--SYRTE link with $\Delta h = 25$~m:
\begin{align}
\delta t_{\rm NR}^{\rm PTB-SYRTE} &= \frac{2g\,\Delta h}{c^3}
\times L \nonumber\\
&= \frac{2\times 9.81\times 25}{(3\times 10^8)^3}
\times 1.415\times 10^9~\text{m} \nonumber\\
&= \frac{490.5}{2.7\times 10^{25}}\times 1.415\times 10^9 \nonumber\\
&= 1.816\times 10^{-23}\times 1.415\times 10^9 \nonumber\\
&= 2.57\times 10^{-14}~\text{s} = 25.7~\text{fs}.
\label{eq:signal_PTB_SYRTE}
\end{align}

In fractional frequency units:
\begin{equation}
\frac{\delta\nu}{\nu}\bigg|_{\rm NR}
= \frac{\delta t_{\rm NR}}{T_{\rm prop}}
= \frac{2.57\times 10^{-14}}{4.72\times 10^{-3}}
= 5.4\times 10^{-12}.
\end{equation}

Wait---let us reconsider.  The nonreciprocal delay is a \emph{static}
offset between one-way travel times in opposite directions.  It
does not show up as a frequency shift but as a timing offset.  The
relevant comparison is between the measured one-way delay $A\to B$
and the delay $B\to A$.

The measurement precision needed is $\sigma_t < \delta t_{\rm NR}/3
= 8.6$~fs for $3\sigma$ detection.

\subsection{Measurement protocol}

\subsubsection{Phase 1: Independent one-way time transfer}

\begin{enumerate}
\item Pre-synchronise clocks at PTB and SYRTE using the standard
      bidirectional technique (this determines the symmetric
      gravitational redshift to $10^{-18}$).
\item Switch to one-way mode: PTB transmits a phase-locked
      optical signal to SYRTE at wavelength $\lambda_1 = 1542.14$~nm.
\item Simultaneously, SYRTE transmits to PTB at
      $\lambda_2 = 1557.36$~nm (wavelength-division multiplexed).
\item Each station records the arrival time of the incoming
      signal using its local optical clock.
\item The one-way delays $\Delta t_{P\to S}$ and $\Delta t_{S\to P}$
      are computed.
\item The nonreciprocal signal is
      $\delta t_{\rm NR} = \Delta t_{P\to S} - \Delta t_{S\to P}$.
\end{enumerate}

\subsubsection{Phase 2: Systematic checks}

\begin{enumerate}
\item Swap wavelengths: PTB at $\lambda_2$, SYRTE at $\lambda_1$.
      $\delta t_{\rm NR}$ must be unchanged (eliminates chromatic
      dispersion).
\item Use time-division multiplexing (alternate $P\to S$ and
      $S\to P$ on same wavelength) as cross-check.
\item Compute and subtract the Sagnac correction:
      $\delta t_{\rm Sagnac} = 2\omega_\oplus A_{\rm enclosed}/c^2$,
      where $A_{\rm enclosed}$ is the fiber-enclosed area projected
      onto the equatorial plane.  For the PTB--SYRTE link at
      latitude $\sim 49^\circ$, $A_{\rm enclosed} \sim 5\times 10^{10}$~m$^2$:
      \begin{equation}
      \delta t_{\rm Sagnac} = \frac{2\times 7.29\times 10^{-5}
      \times 5\times 10^{10}}{9\times 10^{16}}
      = 8.1\times 10^{-11}~\text{s} = 81~\text{ps}.
      \end{equation}
      This is much larger than the \DFD\ signal ($25.7$~fs), so it
      must be subtracted accurately.  The Sagnac correction is known
      from geodetic surveys to $\sim 1\%$, giving a Sagnac residual
      of $\sim 0.8$~ps---still much larger than the signal.
\end{enumerate}

\textbf{Critical issue:} The Sagnac correction uncertainty ($\sim 1$~ps)
exceeds the predicted \DFD\ signal ($25.7$~fs) by a factor of $\sim 40$.
This means the PTB--SYRTE link \emph{as configured} ($\Delta h = 25$~m)
cannot detect the \DFD\ nonreciprocal signal.

\subsection{Optimised configuration: large $\Delta h$}

To overcome the Sagnac background, we need a fiber link with
\emph{large} height difference and \emph{small} enclosed area.

\textbf{Ideal configuration:} A vertical fiber link (zero enclosed
area $\Rightarrow$ zero Sagnac effect) with maximum height difference.

Candidate facilities:
\begin{enumerate}
\item \textbf{Laboratoire Souterrain de Modane (LSM):}
      Underground laboratory at $\Delta h = 1700$~m below the
      Fr\'ejus Road Tunnel, connected to surface by fiber.
      $L_{\rm fiber} \approx 5$~km.
      \begin{align}
      \delta t_{\rm NR}^{\rm LSM} &= \frac{2\times 9.81\times 1700}
      {2.7\times 10^{25}}\times 5\times 10^6 \nonumber\\
      &= 1.234\times 10^{-18}\times 5\times 10^6 = 6.2\times 10^{-12}~\text{s}
      = 6.2~\text{ps}.
      \end{align}
      Sagnac: nearly zero (vertical link encloses minimal area).

\item \textbf{Sanford Underground Research Facility (SURF):}
      $\Delta h = 1478$~m, $L_{\rm fiber} \sim 3$~km.
      \begin{equation}
      \delta t_{\rm NR}^{\rm SURF} = \frac{2\times 9.81\times 1478}
      {2.7\times 10^{25}}\times 3\times 10^6 = 3.2~\text{ps}.
      \end{equation}

\item \textbf{Purpose-built vertical shaft:}
      A 1000~m vertical bore with straight fiber ($L = 1000$~m,
      $\Delta h = 1000$~m):
      \begin{equation}
      \delta t_{\rm NR}^{\rm shaft} = \frac{2\times 9.81\times 1000}
      {2.7\times 10^{25}}\times 10^6 = 0.73~\text{ps}.
      \end{equation}
\end{enumerate}

\subsection{Complete noise budget for LSM configuration}

\begin{table}[t]
\caption{Noise budget for the definitive DFD nonreciprocity test
using the LSM vertical fiber link.  Signal: $6.2$~ps.}
\label{tab:noise_LSM}
\begin{ruledtabular}
\begin{tabular}{lcc}
Source & Single-epoch ($10^4$~s) & $N = 100$ epochs \\
\hline
Clock instability ($10^{-18}$) & $47$~fs & $4.7$~fs \\
Fiber noise (suppressed) & $10$~fs & $1$~fs \\
Temperature drift ($<10$~mK) & $5$~fs & $0.5$~fs \\
Sagnac residual & $<1$~fs & $<0.1$~fs \\
Chromatic dispersion & $<1$~fs & $<0.1$~fs \\
\hline
\textbf{Total noise} & $48$~fs & $4.8$~fs \\
\end{tabular}
\end{ruledtabular}
\end{table}

\textbf{Signal-to-noise ratio:}
\begin{equation}
\text{SNR}(N = 100) = \frac{6.2~\text{ps}}{4.8~\text{fs}} \approx 1300.
\end{equation}

The measurement is decisively feasible.  Even a single epoch
gives SNR $\approx 130$.

\subsection{Integration time requirement}

For a $5\sigma$ detection with a single epoch:
\begin{equation}
\sigma_t < \frac{\delta t_{\rm NR}}{5} = \frac{6.2~\text{ps}}{5}
= 1.24~\text{ps}.
\end{equation}

With current optical clock stability $\sigma_y(\tau)
\sim 10^{-16}/\sqrt{\tau}$, the timing uncertainty after
integration time $\tau$ is:
\begin{equation}
\sigma_t(\tau) = \sigma_y(\tau)\times T_{\rm prop}
= \frac{10^{-16}}{\sqrt{\tau}}\times\frac{L}{c}
= \frac{10^{-16}}{\sqrt{\tau}}\times 1.7\times 10^{-5}~\text{s}.
\end{equation}

For $L = 5$~km ($T_{\rm prop} = 1.7\times 10^{-5}$~s):
$\sigma_t(\tau) = 1.7\times 10^{-21}/\sqrt{\tau}$~s.

Setting $\sigma_t = 1.24\times 10^{-12}$~s:
\begin{equation}
\sqrt{\tau} = \frac{1.7\times 10^{-21}}{1.24\times 10^{-12}}
= 1.37\times 10^{-9},
\end{equation}
giving $\tau \sim 2\times 10^{-18}$~s---i.e., the signal is
detectable essentially instantaneously with optical clock precision.

More realistically, the dominant noise is not clock instability
but \emph{fiber-path} instability.  For a 5~km fiber with active
stabilisation, the residual timing instability is
$\sigma_t^{(\rm fiber)} \sim 10^{-14}~\text{s}/\sqrt{\tau/\text{s}}$.

Setting $\sigma_t^{(\rm fiber)} = 1.24\times 10^{-12}$:
\begin{equation}
\tau_{\rm min} = \left(\frac{10^{-14}}{1.24\times 10^{-12}}\right)^2
= 6.5\times 10^{-5}~\text{s} \approx 65~\mu\text{s}.
\end{equation}

The measurement requires negligible integration time.  In practice,
systematic checks (wavelength swaps, temperature cycles,
Sagnac verification) would extend the campaign to $\sim 1$ week.

\subsection{Falsification criteria}

\begin{enumerate}
\item \textbf{Signal magnitude:} $\delta t_{\rm NR} = 6.2 \pm 0.6$~ps
      (for $\Delta h = 1700 \pm 10$~m, $L = 5000 \pm 50$~m).
      DFD is excluded if $|\delta t_{\rm NR}| < 1.2$~ps at $3\sigma$.

\item \textbf{Sign:} The delay from bottom to top must exceed
      the delay from top to bottom (signal propagating \emph{upward}
      against $\nabla\psi$ is slower in the clock-rate-corrected
      measurement).  Wrong sign excludes DFD.

\item \textbf{Height scaling:} If measurements are made at
      multiple heights (e.g., intermediate levels in the shaft),
      $\delta t_{\rm NR}$ must scale linearly with $\Delta h$.
      Nonlinear scaling excludes DFD.

\item \textbf{Length scaling:} $\delta t_{\rm NR}$ must scale
      linearly with fiber length $L$ (verified by introducing
      fiber spools).

\item \textbf{Two-way null test:} A round-trip measurement
      ($A\to B\to A$) must yield zero nonreciprocal signal to
      within noise.

\item \textbf{Horizontal null test:} A horizontal fiber of the
      same length ($\Delta h \approx 0$) must show zero signal.
\end{enumerate}

%=============================================================================
\section{Discussion}\label{sec:discussion}
%=============================================================================

\subsection{Why the effect has not been detected}

The \DFD\ nonreciprocal timing effect has evaded detection because:

\begin{enumerate}
\item \textbf{Two-way blindness:} The vast majority of precision
      timing experiments (GP-A, Cassini, LLR, radar ranging,
      fiber-link clock comparisons) use two-way or round-trip
      techniques that cancel the nonreciprocal term identically.

\item \textbf{GPS averaging:} GPS position solutions use multiple
      satellites simultaneously, averaging the elevation-dependent
      $\delta t_{\rm NR}$ across different lines of sight.  The
      $\sim 0.2$~ns systematic is comparable to multipath and
      atmospheric residuals.

\item \textbf{Sagnac dominance:} In one-way fiber experiments,
      the Sagnac correction ($\sim 10$--$100$~ps for horizontal
      links) is much larger than the DFD signal ($\sim 1$--$10$~fs
      for small $\Delta h$), masking the effect.

\item \textbf{Vertical links not tested:} No experiment has
      specifically measured one-way timing asymmetry on a vertical
      fiber link with optical-clock precision.
\end{enumerate}

\subsection{Relationship to PPN framework}

The \DFD\ nonreciprocal effect is \emph{invisible} to the
Parameterised Post-Newtonian (PPN) framework~\cite{Will2018}
because PPN analyses the physical metric (governing material
clocks and rulers) rather than separately tracking light
propagation in the optical metric.  In GR, the physical and
optical metrics are identical.  In \DFD, they differ: the
physical metric determines clock rates and particle dynamics,
while the optical metric~\eqref{eq:optical_metric} governs
photon propagation.  The nonreciprocal delay arises from the
\emph{interplay} between clock rates (physical metric) and
light travel times (optical metric)---a degree of freedom that
PPN does not parameterise.

\DFD\ reproduces all PPN parameters exactly:
$\gamma = \beta = 1$, $\xi = \alpha_1 = \alpha_2 = \alpha_3 = 0$,
$\zeta_1 = \zeta_2 = \zeta_3 = \zeta_4 = 0$.  The nonreciprocal
effect is a \emph{post-PPN} prediction.

\subsection{Honest assessment of the Hubble tension connection}

Our detailed calculation (Section~\ref{sec:hubble}) shows that the
simple local $\psi$-well mechanism gives $\delta H_0/H_0 \sim 10^{-4}$,
far too small to explain the $\sim 8\%$ Hubble tension.  The MOND-regime
enhancement increases this by a factor $\sim 100$ but still falls
short by a factor $\sim 600$.

The full resolution of the Hubble tension in \DFD\ requires the
\emph{cosmological} $\psi$-screen, which modifies the global
distance--redshift relation.  This is equivalent to showing that
\DFD\ cosmology naturally produces an effective equation of state
$w_{\rm eff}(z)$ that differs from $\Lambda$CDM in a way that
resolves the tension.  This calculation is an open problem.

We emphasise that the \emph{testable} prediction of this paper is
the GPS/fiber nonreciprocal timing effect, which does not depend
on the Hubble tension connection.

%=============================================================================
\section{Conclusions}\label{sec:conclusions}
%=============================================================================

We have presented the most detailed analysis to date of \DFD\
nonreciprocal timing predictions across all precision timing
systems in the Solar System.  The key results are:

\begin{enumerate}
\item The GPS nonreciprocal correction is $\delta t_{\rm NR}
      = 0.14$--$0.21$~ns, depending on elevation angle
      (Table~\ref{tab:elevation}), with negligible azimuthal
      or seasonal variation.

\item All 15 precision timing experiments in the literature
      are consistent with \DFD\ (Table~\ref{tab:literature}).
      Two-way measurements cancel the effect; one-way
      measurements either lack precision or have not been
      specifically analysed for the nonreciprocal signature.

\item Deep-space one-way ranging corrections range from
      $12$~ns (Mars at closest approach) to $2.1~\mu$s
      (Pluto at encounter), with all major missions using
      two-way ranging that is blind to the effect.

\item The local Milky Way $\psi$-well produces an apparent
      $H_0$ shift of $\sim 10^{-4}$ (MOND-enhanced), too
      small to explain the Hubble tension directly.  The full
      cosmological $\psi$-screen calculation is needed.

\item A vertical fiber link (e.g., LSM, $\Delta h = 1700$~m)
      with optical clocks can detect the predicted
      $6.2$~ps nonreciprocal signal at SNR $> 1000$ in
      $\sim 1$ week of measurements.
\end{enumerate}

The experiment we propose is definitive: it uses existing
infrastructure, the predicted signal exceeds the noise floor
by three orders of magnitude, and six independent falsification
criteria are specified.  If the nonreciprocal delay is detected
with the predicted magnitude, sign, and height dependence, it
would constitute the first evidence for direction-dependent
one-way light propagation in a gravitational field---a
qualitatively new phenomenon invisible to all previous tests
of gravity.

\begin{acknowledgments}
This work is part of the Density Field Dynamics research programme.
Patent provisional application 63/865,279 filed August 16, 2025.
\end{acknowledgments}

%=============================================================================
\begin{thebibliography}{99}

\bibitem{AlcockDFD2025}
G.~T.~Alcock, ``Density Field Dynamics: A Complete Unified Theory,''
v3.2 (2025).

\bibitem{AlcockGPS2026}
G.~T.~Alcock, ``DFD-Corrected GPS: Nonreciprocal Timing, Hubble Bias
Resolution, and Next-Generation Satellite Navigation,''
DFD Research Output, Patent~13 (2026).

\bibitem{Ashby2003}
N.~Ashby, ``Relativity in the Global Positioning System,''
\textit{Living Rev.\ Relativ.}\ \textbf{6}, 1 (2003).

\bibitem{Senior2008}
K.~L.~Senior, J.~R.~Ray, and R.~L.~Beard,
``Characterization of periodic variations in the GPS satellite clocks,''
\textit{GPS Solut.}\ \textbf{12}, 211--225 (2008).

\bibitem{Kouba2009}
J.~Kouba, ``A simplified yaw-attitude model for eclipsing GPS
satellites,'' \textit{GPS Solut.}\ \textbf{13}, 1--12 (2009).

\bibitem{PetitJiang2008}
G.~Petit and Z.~Jiang, ``GPS All in View time transfer for TAI
computation,'' \textit{Metrologia}\ \textbf{45}, S35 (2008).

\bibitem{Vessot1980}
R.~F.~C.~Vessot \textit{et al.}, ``Test of relativistic gravitation
with a space-borne hydrogen maser,''
\textit{Phys.\ Rev.\ Lett.}\ \textbf{45}, 2081 (1980).

\bibitem{HafeleKeating1972a}
J.~C.~Hafele and R.~E.~Keating, ``Around-the-world atomic clocks:
predicted relativistic time gains,''
\textit{Science}\ \textbf{177}, 166--168 (1972).

\bibitem{HafeleKeating1972b}
J.~C.~Hafele and R.~E.~Keating, ``Around-the-world atomic clocks:
observed relativistic time gains,''
\textit{Science}\ \textbf{177}, 168--170 (1972).

\bibitem{Krisher1990}
T.~P.~Krisher, L.~Maleki, G.~F.~Lutes, L.~E.~Primas,
R.~T.~Logan, J.~D.~Anderson, and C.~M.~Will,
``Test of the isotropy of the one-way speed of light using
hydrogen-maser frequency standards,''
\textit{Phys.\ Rev.\ D}\ \textbf{42}, 731 (1990).

\bibitem{Nagel2015}
M.~Nagel \textit{et al.}, ``Direct terrestrial test of Lorentz
symmetry in electrodynamics to $10^{-18}$,''
\textit{Nat.\ Commun.}\ \textbf{6}, 8174 (2015).

\bibitem{Schreiber2004}
K.~U.~Schreiber, T.~Kl\"ugel, and G.~E.~Stedman,
``Earth tide and tilt detection by a ring laser gyroscope,''
\textit{J.\ Geophys.\ Res.}\ \textbf{109}, B06405 (2004).

\bibitem{Lisdat2016}
C.~Lisdat \textit{et al.}, ``A clock network for geodesy and
fundamental science,''
\textit{Nat.\ Commun.}\ \textbf{7}, 12443 (2016).

\bibitem{Delva2017}
P.~Delva \textit{et al.}, ``Test of special relativity using a
fiber network of optical clocks,''
\textit{Phys.\ Rev.\ Lett.}\ \textbf{118}, 221102 (2017).

\bibitem{Titov2011}
O.~Titov, S.~B.~Lambert, and A.-M.~Gontier,
``VLBI measurement of the secular aberration drift,''
\textit{Astron.\ Astrophys.}\ \textbf{529}, A91 (2011).

\bibitem{NANOGrav2023}
G.~Agazie \textit{et al.}\ (NANOGrav Collaboration),
``The NANOGrav 15~yr data set: evidence for a gravitational-wave
background,'' \textit{Astrophys.\ J.\ Lett.}\ \textbf{951}, L8 (2023).

\bibitem{Williams2004}
J.~G.~Williams, S.~G.~Turyshev, and D.~H.~Boggs,
``Progress in lunar laser ranging tests of relativistic gravity,''
\textit{Phys.\ Rev.\ Lett.}\ \textbf{93}, 261101 (2004).

\bibitem{Bertotti2003}
B.~Bertotti, L.~Iess, and P.~Tortora,
``A test of general relativity using radio links with the Cassini
spacecraft,''
\textit{Nature}\ \textbf{425}, 374--376 (2003).

\bibitem{Shapiro1971}
I.~I.~Shapiro \textit{et al.}, ``Fourth test of general
relativity: new radar result,''
\textit{Phys.\ Rev.\ Lett.}\ \textbf{26}, 1132 (1971).

\bibitem{Mueller2010}
H.~M\"uller, A.~Peters, and S.~Chu,
``A precision measurement of the gravitational redshift by the
interference of matter waves,''
\textit{Nature}\ \textbf{463}, 926--929 (2010).

\bibitem{Anderson2002}
J.~D.~Anderson \textit{et al.}, ``Study of the anomalous
acceleration of Pioneer 10 and 11,''
\textit{Phys.\ Rev.\ D}\ \textbf{65}, 082004 (2002).

\bibitem{Turyshev2012}
S.~G.~Turyshev \textit{et al.}, ``Support for the thermal origin
of the Pioneer anomaly,''
\textit{Phys.\ Rev.\ Lett.}\ \textbf{108}, 241101 (2012).

\bibitem{McMillan2017}
P.~J.~McMillan, ``The mass distribution and gravitational
potential of the Milky Way,''
\textit{Mon.\ Not.\ R.\ Astron.\ Soc.}\ \textbf{465}, 76--94 (2017).

\bibitem{Secrest2021}
N.~J.~Secrest \textit{et al.}, ``A test of the cosmological
principle with quasars,''
\textit{Astrophys.\ J.\ Lett.}\ \textbf{908}, L51 (2021).

\bibitem{Will2018}
C.~M.~Will, \textit{Theory and Experiment in Gravitational Physics},
2nd ed.\ (Cambridge Univ.\ Press, 2018).

\bibitem{Riess2022}
A.~G.~Riess \textit{et al.}, ``A comprehensive measurement of the
local value of the Hubble constant,''
\textit{Astrophys.\ J.\ Lett.}\ \textbf{934}, L7 (2022).

\bibitem{Planck2020}
Planck Collaboration, ``Planck 2018 results.\ VI.\ Cosmological
parameters,'' \textit{Astron.\ Astrophys.}\ \textbf{641}, A6 (2020).

\end{thebibliography}

\end{document}
